\documentclass[11pt,a4paper]{ltjsarticle}

% パッケージの読み込み
\usepackage{luatexja} % LuaLaTeXで日本語を扱うための基本パッケージ
\usepackage{graphicx} % 画像挿入用
\usepackage{amsmath,amssymb} % 数式用
\usepackage{hyperref} % リンク・目次用
% \usepackage{k}  % ← 入力途中だったため一時コメントアウト


\title{LuaLaTeX テスト文書}
\author{双子座のAI}
\date{\today}
\begin{document}

\maketitle

\tableofcontents
\clearpage

\section{はじめ}
これLuaLaTeXでコンパイルするためのサンプルテキストです。
LuaLaTeXは、従来のpdfLaTeXやpLaTeXに比べ、Unicodeの扱いが強力で、フォントの設定も柔軟に行えるという特徴があります。

\section{文章のテスト}
通常の日本語文章はもちろん、**太字**や*斜体（※日本語フォントによっては立体になる場合があります）*の表現も可能です。

箇条書きの例：
\begin{itemize}
    \item 最初のアイテム
    \item 2番目のアイテム
    \item 3番目のアイテム
\end{itemize}

\section{数式のテスト}
インライン数式は $E = mc^2$ のように記述します。
独立した数式（ディスプレイスタイル）は以下のように綺麗にレンダリングされます。

\begin{equation}
    \int_{0}^{\infty} e^{-x^2} dx = \frac{\sqrt{\pi}}{2}
\end{equation}

\section{おわりに}
このファイルが無事にPDFとして出力されれば、あなたの環境のLuaLaTeXは正常に動作しています！

\end{document}